\documentclass[10pt,a4paper]{amsart}
\usepackage[margin=19mm]{geometry}
\usepackage{amsmath,amssymb,mathtools}
\usepackage[scr=rsfs,cal=euler]{mathalfa}
\usepackage{xfrac}
\usepackage[unicode,hidelinks,bookmarks=false]{hyperref}
\newcommand{\ie}{i.e.\ }
\newcommand{\cf}{cf.\ }
\newcommand{\resp}{respectively\ }
\newcommand{\Subsection}{\subsection}
\providecommand{\nobreakdash}{\nobreak\hbox{-}}
\setlength{\emergencystretch}{2em}
\multlinegap0pt
\usepackage{bm}
\usepackage{bbm}
\usepackage{dsfont}
\usepackage{xspace}
\usepackage{cancel}
\usepackage{mathtools}
\usepackage{amsthm}
\makeatletter
\@ifundefined{lemma}{\newtheorem{lemma}{Lemma}}{}
\makeatother
\title[Supersonic sonic patch solution for the two-dimensional Eule]{Supersonic sonic patch solution for the two-dimensional Euler equations with a van der Waals equation of state}
\author{Anamika Pandey, T. Raja Sekhar}
\date{Source: arXiv:2512.20484v2 (CC BY 4.0)}
\begin{document}
\maketitle

\noindent\emph{Source and licence.} Supersonic sonic patch solution for the two-dimensional Euler equations with a van der Waals equation of state. Version arXiv:2512.20484v2 (CC BY 4.0), distributed under the Creative Commons Attribution 4.0 International licence, as recorded in the arXiv licence metadata for this exact version and shown on the abstract page https://arxiv.org/abs/2512.20484v2. Full rights notice: licenses/arxiv-2512.20484-CC-BY-4.0.txt.\par\smallskip

\noindent\textit{Editorial scope.} Counted unit: the Lemma whose source label is \texttt{Sol\_3\_l6} in the article (source0.tex lines 1046--1048, source label \texttt{Sol\_3\_l6}), delivered together with the article's own complete original proof of it (source0.tex lines 1049--1159, one \texttt{proof} environment of 111 lines). No mathematics is written, repaired or re-derived: statement and proof are the article's own text line for line. Required context retained uncounted: Sol1\_3 (lines 442--448); Sol1\_4 (lines 450--455); Sol1\_7 (lines 471--479); Sol1\_10 (lines 505--511); Sol\_3\_l1 (lines 603--608); Sol3\_5 (lines 711--718); Sol3\_6 (lines 720--728); Sol\_3\_l3 (lines 894--896); Sol4\_24 (lines 982--989); Sol\_3\_l5 (lines 1000--1002); Sol4\_27 (lines 1009--1011). Every definition, display and label of those lines is retained unchanged. Omitted from this package and not redistributed: 1--441, 449--449, 456--470, 480--504, 512--602, 609--710, 719--719, 729--893, 897--981, 990--999, 1003--1008, 1012--1045, 1160--1630. Nothing inside a retained excerpt is omitted or altered: every retained body is delivered exactly as the article's lines are written, so no omission receipt of any kind accompanies this package. Citations: the retained excerpts carry no citation command at all, so no bibliography entry of the article is transcribed and no reference is redistributed. Reference adaptation: none was needed and none was made; every reference inside the delivered unit resolves inside it, so no unresolved reference or citation is produced. Printed slips kept literal: no printed word, symbol, hypothesis or number of the counted statement or of its proof was altered. Layout and class: the article's own class is \texttt{elsarticle} with packages including amsmath, graphicx, caption, wrapfig, color, soul, enumitem, subcaption, comment, times, amssymb, array, yhmath, latexsym; the wrapper is the lane's pinned \texttt{amsart} editorial wrapper at 10pt on A4, which restates the theorem-like environments the retained text needs and adds the article's own macro definitions that the retained text needs (none), with extra wrapper packages (bm, bbm, dsfont, xspace, cancel, mathtools). The delivered numbering is the unit's own. Assets: no retained span contains a figure environment, a graphics-inclusion command, a PDF-page inclusion, a table or a TikZ picture; no image file is copied into this package, the package declares no asset dependency and no image file is redistributed with it. Licence: version arXiv:2512.20484v2 of this article is distributed under the Creative Commons Attribution 4.0 International licence, as recorded in the arXiv licence metadata for this exact version and shown on the abstract page https://arxiv.org/abs/2512.20484v2. A full rights notice is delivered at licenses/arxiv-2512.20484-CC-BY-4.0.txt. Characters: every retained excerpt is pure ASCII, so the delivered engine, XeLaTeX with its default Latin Modern text font, reports no missing character.
\begin{equation}\label{Sol1_3}
\begin{split}
    {X}|_{\widehat{M'Q'}} &= \hat{a}(\hat{\xi}(z)) =: \hat{a}(z), \\
    { Y}|_{\widehat{M'Q'}} &= \hat{b}(\hat{\xi}(z)) =: \hat{b}(z), \\
   W|_{\widehat{M'Q'}} &=\hat{d}(\hat{\xi}(z)) =: \hat{d}(z) \quad \forall \;z \; \in [0,z_0].
\end{split}
\end{equation}

\begin{equation}\label{Sol1_4}
\begin{split}
    &\hat{a}(z), \hat{b}(z), \hat{d}(z) \in C^1([0,z_0]), \\
    &\hat{m}_0 \leq \hat{a}(z), \hat{b}(z), \hat{d}(z) \leq \hat{M}_0, \quad \forall z \in [0,z_0],
\end{split}
\end{equation}

\begin{equation}\label{Sol1_7}
    \begin{split}
        \begin{cases}
             c(z,t) &= \left[2(1-t^2)\left(\frac{2a}{\tau}-z-\hat\phi(\xi_2)-\frac{K(\gamma\tau+(\tau-b))}{\gamma(\tau-b)^{\gamma+1}} \right) \right]^{1/2},\\
             c_t(z,t) &=-\frac{ct}{(1+\kappa(\tau)-\kappa(\tau) t^2)(1-t^2)},\\ 
             c_z(z,t) &=\frac{t^2-1}{c(1+\kappa(\tau)-\kappa(\tau) t^2)}.
        \end{cases}
    \end{split}
\end{equation}

\begin{equation}\label{Sol1_10}
\begin{split}
    f(t,\tau) &=\frac{1}{(1+\kappa(\tau)-\kappa(\tau)t^2)(1-t^2)^{3/2}}\;, \; g(t,\tau) = -\frac{(1-t^2)^2f}{c},\\
    \mathcal{T}_1 &= \frac{A(\kappa(\tau)(Y-tg)}{f(1-t^2)},\;%,\;\\
    \mathcal{T}_2 = \frac{A(\kappa(\tau))(X-tg)}{f(1-t^2)}.
\end{split}
\end{equation}

\begin{lemma}\label{Sol_3_l1}
     Assume that (\ref{Sol1_4}) holds and $(\overline{X},\overline{ Y})(z,t)$ be a $C^1$ solution of (\ref{Sol1_7}) and (\ref{Sol1_3}) in the domain $\Omega_\epsilon$. Then the solution $(\overline{X},\overline{ Y})(z,t)$ satisfies
     \begin{equation}\label{l11}
       \frac{1}{2}\hat{m}_0 \leq \overline{X}, \overline{Y} \leq 2\hat{M}_0.
     \end{equation}
\end{lemma}

\begin{equation}\label{Sol3_5}
    \begin{split}
        \begin{cases}
          \tilde\partial_+\widetilde X &= \frac{\widetilde X-\widetilde Y}{2t}+H_{11}(\widetilde X-\widetilde Y)+H_{12}t,\\
          \tilde\partial_-\widetilde Y &= -\frac{\widetilde X-\widetilde Y}{2t}+H_{21}(\widetilde X-\widetilde Y)+H_{22}t,
        \end{cases}
    \end{split}
\end{equation}

\begin{equation}\label{Sol3_6}
    \begin{split}
        H_{11}&= -\frac{\sqrt{1-t^2}\widetilde Y}{2c(1+\kappa-\kappa t^2)(1+tg\widetilde Y)},\;\; H_{21}= -\frac{\sqrt{1-t^2}\widetilde X}{2c(1+\kappa-\kappa t^2)(1-tg\widetilde X)},\\
        H_{12} &=\frac{(\widetilde X-\widetilde Y)(1+2\kappa-\kappa t^2)}{2(1+tg\widetilde Y)(1-t^2)(1+\kappa-\kappa t^2)}-\frac{f\sqrt{1-t^2}}{1+tg\widetilde Y}\left(\frac{2t\sqrt{1-t^2}\widetilde X\widetilde Y}{c}-\left(\tau\kappa'(\tau)-1-2\kappa(\tau)+2t^2\kappa(\tau)\right)\widetilde X \right),\\
        H_{22} &=-\frac{(\widetilde X-\widetilde Y)(1+2\kappa-\kappa t^2)}{2(1-tg\widetilde X)(1-t^2)(1+\kappa-\kappa t^2)}+\frac{f\sqrt{1-t^2}}{1-tg\widetilde X}\left(\frac{2t\sqrt{1-t^2}\widetilde X\widetilde Y}{c}-\left(\tau\kappa'(\tau)-1-2\kappa(\tau)+2t^2\kappa(\tau)\right)\widetilde Y \right).
        %T_3&=\sqrt{1+\kappa-\kappa t^2}-\frac{t\sqrt{1-t^2}\widetilde Y}{c\sqrt{1+\kappa-\kappa t^2}},\;\;
        %T_4=\sqrt{1+\kappa-\kappa t^2}+\frac{t\sqrt{1-t^2}\widetilde X}{c\sqrt{1+\kappa-\kappa t^2}}
    \end{split}
\end{equation}

\begin{lemma}\label{Sol_3_l3}
    Let $\nu\in(0,1)$ be fixed. Then the quantities $t^{\nu}|\widetilde  {X}_{z}|$ and $t^{\nu}| \widetilde Y_{z}|$ remain uniformly bounded throughout the entire domain $\Omega$ up to the degenerate line $t=0$.
\end{lemma}

    \begin{equation}\label{Sol4_24}
      \begin{split}
          \begin{cases}
               \tilde \partial_+ \widetilde W&=(H_{11}-H_{21})\widetilde W+\frac{1}{2}(H_{12}-H_{22})-\frac{1}{2}t^{1-\nu}\widetilde \Theta,\\
               \tilde \partial_- \widetilde W&=(H_{11}-H_{21})\widetilde W+\frac{1}{2}(H_{12}-H_{22})-\frac{1}{2}t^{1-\nu}\widetilde \Xi.
          \end{cases}
      \end{split}
    \end{equation}

\begin{lemma}\label{Sol_3_l5}
    The functions $( X, Y,W)$ are uniformly $C^{1/3}$ continuous throughout the entire domain $\Omega$, including the degenerate line $\widehat{M'N'}$.
\end{lemma}

    \begin{equation}\label{Sol4_27}
        K_1t_m\leq |z_2-z_1|^{1/3}\leq K_2t_m,
    \end{equation}

\begin{lemma}\label{Sol_3_l6}
    The functions $( X, Y)(z,t)$ are uniformly $C^{1-\nu}$ continuous and $W(z,t)$ is uniformly $C^{(2-\nu)/3}$ continuous for any $\nu\in(0,1)$ throughout the domain $\Omega$ including the degenerate line $\widehat{M'N'}$.
\end{lemma}

 \begin{proof}
     It is sufficient to prove that $(\widetilde X, \widetilde Y)$ are uniformly $C^{1-\nu}$ continuous and that $W$ is uniformly $C^{(2-\nu)/3}$ continuous for any $\nu\in(0,1)$. Let $M^{\ast}$ denote a uniform positive constant such that, for $i,j=1,2$,
     \begin{equation}\label{Sol4_32}
         \left | \frac{cf\widetilde X}{1-tg\widetilde X}\right|, \left | \frac{cf\widetilde Y}{1+tg\widetilde Y}\right|, |2tH_{ij}+1|, |\widetilde W|, |H_{ij\widetilde X}|, |H_{ij\widetilde Y}|,|H_{ij\widetilde Y}c_z|\leq M^*.
     \end{equation}
     The inequalities above follow from Lemma \ref{Sol_3_l1} together with the explicit formulas for $f$, $g$ and $H_{ij}$ given in (\ref{Sol3_6}) and (\ref{Sol1_10}). Furthermore, by Lemma \ref{Sol_3_l3}, there exists a uniform positive constant $a_{2}$, depending on the parameter $\nu$, such that
     \begin{equation}\label{Sol4_33}
         |\widetilde\Xi|, |\widetilde\Theta|, t^\nu|\widetilde X_z|, t^\nu|\widetilde Y_z|\leq a_2.
     \end{equation}
     We begin by refining the uniform regularity of $\widetilde W$. Recalling the first equation satisfied by $\widetilde W$ in (\ref{Sol4_24}), we obtain
     \begin{equation}\label{Sol4_34}
         \widetilde W_t=-\frac{cft\widetilde Y}{2(1+tg\widetilde Y)}(\widetilde X_z-\widetilde Y_z) +(H_{11}-H_{21})\widetilde W+\frac{1}{2}(H_{12}-H_{22})-\frac{1}{2}t^{1-\nu}\widetilde \Theta.
     \end{equation}
     Integrating (\ref{Sol4_34}) along the segment from $(z_{i},0)$ to $(z_{i},t_{m})$ for $i=1,2$ yields
     \begin{equation}\label{Sol4_35}
          \widetilde W(z_i,0)=\widetilde W(z_i,t_m)-\int_0^t\left(-\frac{cft\widetilde Y}{2(1+tg\widetilde Y)}(\widetilde X_z-\widetilde Y_z) +(H_{11}-H_{21})\widetilde W+\frac{1}{2}(H_{12}-H_{22})-\frac{1}{2}t^{1-\nu}\widetilde \Theta \right)dt.
     \end{equation}
     By using (\ref{Sol4_35}) we obtain
     \begin{equation}\label{Sol4_36}
          \begin{split}
              |\widetilde W(z_1,0)-\widetilde W(z_2,0)|&\leq |\widetilde W(z_1,t_m)-\widetilde W(z_2,t_m)| \\&+\int_0^{t_m}\left(\left|\frac{cf\widetilde Y (\widetilde X_z t^\nu-\widetilde Y_z t^\nu)}{2(1+tg\widetilde Y)}(z_1,t) \right|+\left|\frac{cf\widetilde Y (\widetilde X_z t^\nu-\widetilde Y_z t^\nu)}{2(1+tg\widetilde Y)}(z_2,t) \right|\right)t^{1-\nu}dt\\
              &+\int_0^{t_m} |(H_{11}-H_{21})(z_1,t)|\cdot |\widetilde W(z_1,t)-\widetilde W(z_2,t)| dt \\
              &+ \int_0^{t_m}  |(H_{11}-H_{21})(z_1,t)-(H_{11}-H_{21})(z_2,t)|\cdot |\widetilde W(z_2,t)|dt\\
              &+\int_0^{t_m} \frac{1}{2}|(H_{12}-H_{22})(z_1,t)-(H_{12}-H_{22})(z_2,t)|dt+\int_0^{t_m}\frac{1}{2}\left(|\widetilde \Theta(z_1,t)|+|\widetilde \Theta(z_2,t)| \right)t^{1-\nu} dt.
          \end{split}
     \end{equation}
     Consider the leading term on the right hand side of (\ref{Sol4_36}). Using the estimates provided in (\ref{Sol4_33})  and (\ref{Sol4_27}), we deduce that
     \begin{equation}\label{Sol4_37}
         \begin{split}
             |\widetilde W(z_1,t_m)-\widetilde W(z_2,t_m)|&\leq \frac{1}{t_m}\left( \left|\widetilde X(z_1,t_m)-\widetilde X(z_2,t_m)\right|+\left|\widetilde Y(z_1,t_m)-\widetilde Y(z_2,t_m)\right|\right)\\
             &\leq  t_{m}^{-1-\nu}\left(\left|t_m^\nu\widetilde X(z',t_m)\right|\cdot|z_1-z_2|+ \left|t_m^\nu\widetilde Y(z'',t_m)\right|\cdot|z_1-z_2|\right)\\
             &\leq 2a_2t_{m}^{-1-\nu}|z_1-z_2|\leq 2a_2 \;K_2^{1+\nu}|z_1-z_2|^\frac{2-\nu}{3}.
         \end{split}
     \end{equation}
     Furthermore, since the functions $(\widetilde X,\widetilde Y,\widetilde W)(z,t)$ possess uniform $C^{1/3}$ continuity as established in Lemma \ref{Sol_3_l5}, it follows that
     \begin{equation}\label{Sol4_38}
       |\widetilde X(z_1,t)-\widetilde X(z_2,t)|, |\widetilde Y(z_1,t)-\widetilde Y(z_2,t)|, |\widetilde W(z_1,t)-\widetilde W(z_2,t)|\leq a_3|z_1-z_2|^{1/3},
     \end{equation}
     for a certain uniform positive constant $a_3$. Consequently, combining (\ref{Sol4_32}) and (\ref{Sol4_38}), we conclude that for $i=1,2$,
     \begin{equation*}\label{Sol4_39}
         \begin{split}
             \left| (H_{1i}-H_{2i})(z_1,t)-(H_{1i}-H_{2i})(z_2,t)\right| &\leq |H_{1i\widetilde X}-H_{2i\widetilde X}|\cdot|\widetilde X(z_1,t)-\widetilde X(z_2,t)|+ |H_{1i\widetilde Y}-H_{2i\widetilde Y}|\\& \cdot|\widetilde Y(z_1,t)-\widetilde Y(z_2,t)|+|H_{1ic}-H_{2ic}|\cdot c_z\cdot |z_1-z_2|\\
             &\leq 4M^*a_3|z_1-z_2|^{1/3}+2M^*|z_1-z_2|.
         \end{split}
     \end{equation*}
     We now substitute (\ref{Sol4_37}) and (\ref{Sol4_38}) into (\ref{Sol4_36}) and then apply the bounds given in  (\ref{Sol4_27}), (\ref{Sol4_32}) and (\ref{Sol4_33}) to obtain
     \begin{equation}\label{Sol4_40}
         \begin{split}
             |\widetilde W(z_1,0)-\widetilde W(z_2,0)|\leq &2a_2 \;K_2^{1+\nu}|z_1-z_2|^\frac{2-\nu}{3}+\int_0^{t_m} a_2(2M^*+1)t^{1-\nu} dt\\
             &+ \int_0^{t_m} 2M^* a_3|z_1-z_2|^{1/3} dt+\int_0^{t_m} \left(M^*+\frac{1}{2}\right)\left(4M^*a_3|z_1-z_2|^{1/3}+2M^*|z_1-z_2|\right)dt\\
             \leq & 2a_2 \;K_2^{1+\nu}|z_1-z_2|^\frac{2-\nu}{3}+\frac{2M^*+1}{2-\nu}a_2t_m^{2-
             \nu}+2M^*a_3|z_1-z_2|^{1/3}t_m\\
             &+ (M^*+1)\left(4M^*a_3|z_1-z_2|^{1/3}+2M^*|z_1-z_2| \right)t_m\\
             \leq & a_4 |z_1-z_2|^{\frac{2-\nu}{3}},
         \end{split}
     \end{equation}
     for a uniform positive constant $a_4$ that depends on $\nu$. In view of (\ref{Sol4_40}), the same argument employed in Lemma \ref{Sol_3_l5} can be repeated to conclude that the function $\widetilde W$ is uniformly $C^{(2-\nu)/3}$ continuous throughout the entire domain $\Omega$.

     For the functions $\widetilde X$ and $\widetilde Y$, we return to (\ref{Sol3_5}) to obtain
     \begin{equation}\label{Sol4_41}
         \begin{split}
             \begin{cases}
                 \widetilde X_t =-\frac{cft^2\widetilde Y}{1+tg\widetilde Y}\widetilde X_z+(2tH_{11}+1)\widetilde W+H_{12}t,\\
                 \widetilde Y_t =\frac{cft^2\widetilde X}{1-tg\widetilde X}\widetilde Y+(2tH_{21}-1)\widetilde W+H_{22}t.
             \end{cases}
         \end{split}
     \end{equation}
     Integrating the equation satisfied by $\widetilde X$ in (\ref{Sol4_41}) along the path from $(z_{i},0)$ to $(z_{i},t_{m})$ for $i=1,2$, we obtain
     \begin{equation*}\label{Sol4_42}
         \widetilde X(z_i,0)=\widetilde X(z_i,t_m)-\int_0^{t_m} \left(-\frac{cft^2\widetilde Y}{1+tg\widetilde Y}+(2tH_{11}+1)\widetilde W+H_{12}t\right)(z_i,t)dt,
     \end{equation*}
     from which it follows that
     \begin{equation}\label{Sol4_43}
         \begin{split}
             |\widetilde X(z_1,0)-\widetilde X(z_2,0)|\leq & |\widetilde X(z_1,t_m)-\widetilde X(z_2,t_m)|+\int_0^{t_m} |H_{12}(z_1,t)-H_{12}(z_2,t)|tdt\\
              +\int_0^{t_m}&\left( \left|\frac{cf\widetilde Y}{1+tg\widetilde Y}(z_1,t)\right| \cdot|t^\nu\widetilde X_z(z_1,t)|+ \left|\frac{cf\widetilde Y}{1+tg\widetilde Y}(z_2,t)\right| \cdot|t^\nu\widetilde X_z(z_2,t)|\right)t^{2-\nu}dt\\
              +\int_0^{t_m}&\left( |2tH_{11}(z_1,t)+1|\cdot|\widetilde W(z_1,t)-\widetilde W(z_2,t)|+2t|\widetilde W|+2t|\widetilde W|\cdot|H_{11}(z_1,t)-H_{11}(z_2,t)| \right)dt.
         \end{split}
     \end{equation}
     By making use of the bounds in (\ref{Sol4_32}) and (\ref{Sol4_33}), we conclude that
     \begin{equation}\label{Sol4_44}
         \begin{split}
             |\widetilde X(z_1,t_m)-\widetilde X(z_2,t_m)|&=|\widetilde X_z(z',t_m)|\cdot|z_1-z_2|\leq a_2t_m^{-\nu} |z_1-z_2|\leq a_2{K_2}^\nu |z_1-z_2|^{1-\nu}.
         \end{split}
     \end{equation}
     Also, we have
     \begin{equation*}\label{Sol4_45}
         \left|\frac{cf\widetilde Y}{1+tg\widetilde Y}(z_1,t)\right| \cdot|t^\nu\widetilde X_z(z_1,t)|+ \left|\frac{cf\widetilde Y}{1+tg\widetilde Y}(z_2,t)\right| \cdot|t^\nu\widetilde X_z(z_2,t)|\leq 2M^*a_2.
     \end{equation*}
     Moreover, (\ref{Sol4_32}) implies that for $i=1,2$,
     \begin{equation*}\label{Sol4_46}
         \begin{split}
             |H_{1i}(z_1,t)-H_{1i}(z_2,t)|\leq& |H_{1i\widetilde X}|\cdot |\widetilde X(z_1,t)-\widetilde X(z_2,t)|+|H_{1i\widetilde Y}|\cdot |\widetilde Y(z_1,t)-\widetilde Y(z_2,t)|\\
             &+|H_{1i\tilde c}|\cdot |\tilde c(z_1,t)-\tilde c(z_2,t)|\\
             \leq &\left(2a_2t^{-\nu}+1 \right)M^*|z_1-z_2|.
         \end{split}
     \end{equation*}
     As a consequence of the uniform $C^{(2-\nu)/3}$ continuity of $\widetilde W$, it follows that
     \begin{equation}\label{Sol4_47}
         |\widetilde W(z_1,t)-\widetilde W(z_2,t)|\leq a_5|z_1-z_2|^{\frac{2-\nu}{3}},
     \end{equation}
     for a uniform positive constant $a_5$ depending on $\nu$. Employing \eqref{Sol4_44}-\eqref{Sol4_47} into \eqref{Sol4_43} and then using \eqref{Sol4_27} and \eqref{Sol4_32}, we conclude that
     \begin{equation}\label{Sol4_48}
         \begin{split}
              |\widetilde X(z_1,0)-\widetilde X(z_2,0)|&\leq a_2K_2^\nu|z_1-z_2|^{1-\nu}+2M^*a_2\int_0^{t_m} t^{2-\nu} dt\\
              &+M^*\int_0^{t_m} \left(a_5|z_1-z_2|^{\frac{2-\nu}{3}}+t(2M^*+1)(2a_2t^{-\nu}+1)|z_1-z_2| \right)dt\\
              &= a_2K_2^\nu|z_1-z_2|^{1-\nu}+\frac{2M^*a_2}{3-\nu}t_m^{3-\nu} +M^*a_5|z_1-z_2|^{\frac{2-\nu}{3}}t_m+(2M^*+1)(2a_2+1)M^*|z_1-z_2|\\&\leq a_6|z_1-z_2|^{1-\nu},
         \end{split}
     \end{equation}
     for a uniform positive constant $a_6$. Based on (\ref{Sol4_48}), arguments analogous to those used in Lemma \ref{Sol_3_l5} can be applied to conclude that the function $\widetilde X$ is uniformly $C^{1-\nu}$ continuous over the entire domain $\Omega$. By the same reasoning, the function $\widetilde Y$ also enjoys uniform $C^{1-\nu}$ continuity. This completes the proof of the lemma.
 \end{proof}

\end{document}
